% GENERATED FILE. Edit canonical lessons, then run npm run generate:books.
% canonical-source-hash: fnv1a64:58f749e8bb10c55a
% canonical-lessons: PA-C209-nila-return, PA-C209-nili, PA-C209-nili-review, PA-C209-blue-shoe, PA-C209-payoff

\chapter{This Is a Blue Shoe}
\label{ch:pa-a2-nili}
\hlchaptermodality{209}

\begin{chapteropening}
\textbf{By the end of this chapter:} \emph{I can recognise and use one bounded feminine singular adjective form.}

Recognise and produce one bounded blue-shoe frame through listening, speaking, reading, and writing.
\end{chapteropening}

\section[\texorpdfstring{\pa{ਨੀਲਾ} · \pa{ਜੁੱਤੀ} (nila; jutti)}{ਨੀਲਾ · ਜੁੱਤੀ (nila; jutti)}]{Return to blue and a shoe}

\label{lesson:PA-C209-nila-return}



\begin{quote}
\emph{Hear:} \emph{nila; jutti}; YOU SAY: \textbf{\pa{ਨੀਲਾ} · \pa{ਜੁੱਤੀ}}
\end{quote}

\begin{grammarlens}[title={Grammar Lens}]
Keep the familiar colour word and the familiar noun ready. The next lesson will change one ending to own one exact agreeing form; this opener does not reveal that form.
\end{grammarlens}

\subsection*{Your turn}
\begin{itemize}
\raggedright
  \item \emph{Hear:} \emph{blue}; YOU SAY: \textbf{\pa{ਨੀਲਾ}}
  \item \emph{Hear:} \emph{a shoe}; YOU SAY: \textbf{\pa{ਜੁੱਤੀ}}
  \item Say the two familiar items once, then stop.
\end{itemize}

\begin{tcolorbox}[breakable,colback=teal!4,colframe=teal!35!black,title={Before you move on}]
Produce “blue” and “a shoe” from meaning. This opener owns no new form.
\end{tcolorbox}
\section[\texorpdfstring{\pa{ਨੀਲੀ} (nili)}{ਨੀਲੀ (nili)}]{One bounded feminine adjective form}

\label{lesson:PA-C209-nili}



\begin{quote}
Retrieve familiar \textbf{\pa{ਨੀਲਾ}}, “blue,” and the familiar rule that a variable adjective agrees with a feminine noun.
\end{quote}

\begin{grammarlens}[title={Grammar Lens: direct owner}]
\begin{quote}
\textbf{\pa{ਨੀਲਾ}} — \emph{nīlā} — blue, masculine singular direct
\end{quote}

\begin{quote}
\textbf{\pa{ਨੀਲੀ}} — \emph{nīlī} — blue, feminine singular direct
\end{quote}

For this one form, change the final \textbf{-\pa{ਾ}} ending of familiar \textbf{\pa{ਨੀਲਾ}} to feminine singular \textbf{-\pa{ੀ}}. The result is \textbf{\pa{ਨੀਲੀ}}. Keep the stem and meaning steady; only the ending changes.
\end{grammarlens}

\begin{grammarlens}[title={Grammar Lens: source boundary}]
Punjabi University Patiala's reference grammar says that inflectable adjectives agree with their head noun in gender, number, and case. Its noun-phrase examples contrast masculine \textbf{\pa{ਚੰਗਾ} \pa{ਮੁੰਡਾ}} with feminine \textbf{\pa{ਚੰਗੀ} \pa{ਕੁੜੀ}}. Source: \href{https://pt.learnpunjabi.org/assets/a\%20reference\%20grammar\_final.pdf}{Punjabi University reference grammar}.

That source supports this one feminine singular direct ending change. This chapter does not teach a full agreement paradigm, plural or oblique forms, a new colour, or a general contrast between attributive and predicative uses.
\end{grammarlens}

\subsection*{Your turn}
\begin{itemize}
\raggedright
  \item \emph{Hear:} \emph{nili, blue for one feminine noun}; YOU SAY: \textbf{\pa{ਨੀਲੀ}}
  \item \emph{Read it:} \textbf{\pa{ਨੀਲਾ} · \pa{ਨੀਲੀ}}
  \item \emph{Write it:} copy \textbf{\pa{ਨੀਲੀ}} once, checking the final \textbf{-\pa{ੀ}}
  \item \emph{Check:} familiar \textbf{\pa{ਨੀਲਾ}} ends in \textbf{-\pa{ਾ}}; exact feminine singular \textbf{\pa{ਨੀਲੀ}} ends in \textbf{-\pa{ੀ}}
\end{itemize}

\begin{tcolorbox}[breakable,colback=teal!4,colframe=teal!35!black,title={Before you move on}]
Produce familiar “blue,” then its one feminine singular direct form.
\end{tcolorbox}
\section[\texorpdfstring{\pa{ਨੀਲੀ} (nili)}{ਨੀਲੀ (nili)}]{Retrieve the new feminine form on its own}

\label{lesson:PA-C209-nili-review}



\begin{quote}
\emph{Hear:} \emph{nili}; YOU SAY: \textbf{\pa{ਨੀਲੀ}}
\end{quote}

\subsection*{Your turn}
\begin{itemize}
\raggedright
  \item \emph{Hear:} \emph{blue, feminine singular direct}; YOU SAY: \textbf{\pa{ਨੀਲੀ}}
  \item \emph{Read it:} \textbf{\pa{ਨੀਲੀ}}
  \item \emph{Write it:} cover and write \textbf{\pa{ਨੀਲੀ}}
\end{itemize}

\begin{grammarlens}[title={Grammar Lens}]
Keep one exact form and one bounded agreement job here. The familiar noun and sentence frame return only after this separate retrieval.
\end{grammarlens}

\begin{tcolorbox}[breakable,colback=teal!4,colframe=teal!35!black,title={Before you move on}]
Say and then write the exact form once from memory.
\end{tcolorbox}
\section[\texorpdfstring{\pa{ਇਹ} \pa{ਨੀਲੀ} \pa{ਜੁੱਤੀ} \pa{ਹੈ} (eh nili jutti hai)}{ਇਹ ਨੀਲੀ ਜੁੱਤੀ ਹੈ (eh nili jutti hai)}]{Put the new form before one familiar feminine noun}

\label{lesson:PA-C209-blue-shoe}



\begin{quote}
Retrieve the new feminine colour form, familiar “shoe,” familiar “this,” and familiar “is.”
\end{quote}

\begin{grammarlens}[title={Grammar Lens}]
\begin{quote}
\textbf{\pa{ਇਹ} \pa{ਨੀਲੀ} \pa{ਜੁੱਤੀ} \pa{ਹੈ}} — \emph{eh nīlī juttī hai} — This is a blue shoe.
\end{quote}

For this one frame, begin with familiar \textbf{\pa{ਇਹ}}, place new \textbf{\pa{ਨੀਲੀ}} before the familiar feminine noun \textbf{\pa{ਜੁੱਤੀ}}, and keep familiar \textbf{\pa{ਹੈ}} at the end. The adjective is feminine singular direct because it agrees with this one noun.
\end{grammarlens}

\subsection*{Your turn}
\begin{itemize}
\raggedright
  \item \emph{Hear:} \emph{eh nili jutti hai}; YOU SAY: \textbf{\pa{ਇਹ} \pa{ਨੀਲੀ} \pa{ਜੁੱਤੀ} \pa{ਹੈ}}
  \item \emph{Read it:} \textbf{\pa{ਇਹ} \pa{ਨੀਲੀ} \pa{ਜੁੱਤੀ} \pa{ਹੈ}}
  \item \emph{Write it:} cover and write \textbf{\pa{ਇਹ} \pa{ਨੀਲੀ} \pa{ਜੁੱਤੀ} \pa{ਹੈ}}
  \item \emph{Check:} \textbf{\pa{ਨੀਲੀ}} comes before \textbf{\pa{ਜੁੱਤੀ}}; \textbf{\pa{ਹੈ}} waits at the end
\end{itemize}

\begin{tcolorbox}[breakable,colback=teal!4,colframe=teal!35!black,title={Before you move on}]
Produce the complete “This is a blue shoe” frame.
\end{tcolorbox}
\section[\texorpdfstring{\pa{ਇਹ} \pa{ਨੀਲੀ} \pa{ਜੁੱਤੀ} \pa{ਹੈ} (eh nili jutti hai)}{ਇਹ ਨੀਲੀ ਜੁੱਤੀ ਹੈ (eh nili jutti hai)}]{A bounded blue-shoe payoff}

\label{lesson:PA-C209-payoff}



\begin{quote}
Retrieve the new feminine colour form alone before rebuilding its familiar frame.
\end{quote}

\subsection*{Your turn: four skills}
\begin{itemize}
\raggedright
  \item \emph{Hear:} \emph{eh nili jutti hai}; YOU SAY: \textbf{\pa{ਇਹ} \pa{ਨੀਲੀ} \pa{ਜੁੱਤੀ} \pa{ਹੈ}}
  \item \emph{Read it:} \textbf{\pa{ਇਹ} \pa{ਨੀਲੀ} \pa{ਜੁੱਤੀ} \pa{ਹੈ}}
  \item \emph{Write it:} cover the line and write \textbf{\pa{ਇਹ} \pa{ਨੀਲੀ} \pa{ਜੁੱਤੀ} \pa{ਹੈ}}
\end{itemize}

\begin{tcolorbox}[breakable,colback=teal!4,colframe=teal!35!black,title={Before you move on}]
Say the new form, then produce the complete “This is a blue shoe” frame.

This is the thirty-second prerequisite tranche exposed by Punjabi A2 mock drafting. It does not unblock Punjabi A2 mock 1 by itself. It is not an A2 mock, not A2 readiness or attainment, and not external or human validation.
\end{tcolorbox}
